% From mitthesis package
% Version: 1.01, 2023/06/19
% Documentation: https://ctan.org/pkg/mitthesis
%
% The abstract environment creates all the required headers and footnote. 
% You only need to add the text of the abstract itself.
%
% Approximately 500 words or less; try not to use formulas or special characters
% If you don't want an initial indentation, do \noindent at the start of the abstract

Emotions are an important part of our socio-cognitive toolkit, they help us prioritise our attention resources and guide our decision-making processes, resulting in more 
apt responses. A great deal of success in social contexts could be attributed to our ability to use our emotions to better fit our responses to the environment. Emotion 
regulation is an important aspect of affective behaviours, and has an important role in the investigation of psychopathological states. It is our inherent ability to 
modulate our emotional responses to better fit the needs of the environment. Affective flexibility is the ability to transition across various emotional states while 
engaging in a task. Consequently, having a mechanistic understanding of these phenomena is crucial to understanding emotional dysregulation as seen in various 
neuropsychiatric disorders. A major challenge in studying these phenomena has been the absence of a robust measure of the emotional flux of an individual. Traditional 
self-report questionnaires are susceptible to being biased by the subject's own percept of their mental state, and lack the granularity to capture quite many features 
necessary to study these phenomena. We propose a novel measure derived from the experience sampling method and actigraphy, allowing for the capture of the emotional flux 
of an individual through their motion on a two-dimensional emotional space. Furthermore, we demonstrate that this measure can be used to study affective flexibility, 
effectively capturing individual variability in emotional regulation. 



%\footnote{Text from Holman (1876): \doi{10.2307/25138434}.}  
